\section{Es geht los}

\begin{frame}{Worum geht es?}
    \begin{itemize}
    \item Haben eine Folge $\mathbb{X} = (X_n)_{n \in \N}$ von unabhängigen zentrierten Zufallsvariablen mit Varianz 1 (mit gewissen Momentannahmen).  
    \item Interessiert an folgenden homogenen Summen: \[
    Q_d(f^{(N)}, N, \mathbb{X}) = \sum_{1 \leq i_1, \ldots, i_d \leq  N} f^{(N)}(i_1, \ldots, i_d) X_{i_1} \cdots X_{i_d}\]
    für symmetrische $f: \{1, \ldots, N\}^d \to \R$; $f^{(N)}(i_1, \ldots, i_d) = 0$ falls $i_j = i_k$ für $j \neq k$.
    \pause
    \item Was passiert in Verteilung, falls $N \to \infty$? 
    \item Im Prinzip einiges möglich: $\chi^2$-Verteilung, Normalverteilung, womöglich noch mehr...
    \end{itemize}
\end{frame}

\begin{frame}{Worum geht es?}
    \begin{itemize}
        \item Es genügt jedoch, ein wesentlich einfacheres Objekt zu betrachten: \[
        Q_d(f^{(N)}, N, \mathbb{G}) = \sum_{1 \leq i_1, \ldots, i_d \leq  N} f^{(N)}(i_1, \ldots, i_d) G_{i_1} \cdots G_{i_d}\]
        für $G_1, G_2, \ldots$ unabhängige standardnormalverteilte Zufallsvariablen.
        \pause
        \item Verhält sich asymptotisch unter geeigneten Annahmen genauso.
    \end{itemize}
\end{frame}

\begin{frame}{Warum ist das interessant?}
    \begin{itemize}
        \item Grundsätzlich möglich, bestimmte Bilinearformen abzudecken, und potentielle Anwendungen auf Spektralstatistiken
        \pause
        \item Wir können somit mithilfe des Wiener Chaos andere Chaos-Objekte besser verstehen. Wiener Chaos selbst gut verstanden.
%         \pause
% \begin{satzblockname}{Poisson-Chaos}
%     Für den kompensierten Poisson-Prozess $M_n := P_{\lambda,t} - \lambda t$ mit Intensität $\lambda$ definieren wir \[
%     I_n(f_n) = n! \int_0^\infty \int_0^{t_n} \cdots \int_0^{t_2} f_n(t_1, \ldots, t_n) dM_{t_1} \cdots dM_{t_n}.
%     \]
%     Der von $\{I_k(f_1 \otimes \cdots \otimes f_k) : f_i \in \mathcal{C}_c^\infty(\R_+)\}$ erzeugte Vektorraum entspricht den Polynomen vom Grad $\leq n$ 
%     in $I_1(f), f \in C_c^\infty(\R_+)$.
% \end{satzblockname}
        % \item Zum Beispiel wichtig für $U$-Statistiken-Theorie.
        % \item $U$-Statistiken: \[
        % f(X_1, \ldots, X_n) = \frac{1}{\binom{n}{d}} \sum_{1 \leq i_1 \neq \ldots \neq i_d \leq n} f(X_{i_1} \cdots X_{i_d})\]
        % \item Elementares Beispiel für $d=2$ wäre \[
        % 2 \cdot \binom{n}{2}^{-1}\sum_{1 \leq i_1 < i_2\leq N} X_{i_1} X_{i_2}
        % \]
        % \pause
        % \item Falls lineare Terme dabei, asymptotisch eine Normalverteilung.
        % \item In unserem Fall: "degenerierte" $U$-Statistik, da $f^{(N)}$ keine linearen Terme enthält.
    \end{itemize}
\end{frame}
% \begin{frame}{Wo ist das einzuordnen?}
% \begin{itemize}
%     \item Für $U$-Statistiken der Form \[
%     \sum_{1 \leq i_1 < \ldots < i_d \leq N} f(X_1, \ldots X_{i_d}); X_i \text{ i.i.d.}\] 
%     und $f$ mit $\E[f(X_1, x_2, \ldots, x_n)] = 0$ f.s. ist Asymptotik bereits gut verstanden (präsentiert in \cite{ghs})
%     \pause
%     \item \textit{Unsere} Resultate (aus \cite{homsums}) liefern explizite bounds für Verteilungsabstände, genauer für\[
%     d_{\mathcal{H}}(F, Z) = \sup_{h \in \mathcal{H}} \left|\E[h(F)]- \E[h(Z)]\right|
%     \] und wir betrachten allgemeinere Klassen von $f^{(N)}$ und $X_n$.
% \end{itemize}
% \end{frame}
\begin{frame}{Fahrplan}
    \begin{itemize}
        \item Präsentation des Hauptresultates 
        \item Kontraktionen, Einflussindizes und die Beweisstrategie 
        \item Beweis des Hauptresultates
    \end{itemize}
\end{frame}